\section*{Limitations}
\label{limitations}
The evaluations presented in this study are met with certain limitations due to inherent disparities in the context window of current \emph{open-source} and \emph{proprietary} LLMs and LMMs. 
There is a significant difference in context window length between models such as GPT-4 Turbo and Claude-3, which can rank all papers in a paper cluster, and models such as InternVL-Chat-V1.1 and QwenVL, which are restricted to handling only two to eight papers in a single prompt. 
This discrepancy may lead to an ``unfair'' comparison of their capabilities. 
Future work could focus on standardizing or extending the context windows in LMMs to mitigate this issue. 

Furthermore, as discussed in Section~\ref{main_result}, prompting an LMM with a set of possible reference papers may be suboptimal due to the challenges models face in ranking a large number of papers. An alternative approach could involve assessing the relevance of each paper individually by encoding the paper into a textual embedding, then comparing it with the textual embedding with of the visual context question combined with the image representation of the figure. This method could potentially alleviate the challenges of requiring an LMM to sift through a large set of possible reference papers and would be an interesting area for future research.

Additionally, our approach to ranking papers for certain models, in particular BM25 and Contriever, involves using GPT-4o's textual descriptions of images rather than its direct image embedding, which might not accurately capture the nuances of scientific images. 
Current image embedding models such as LLaVA \cite{liu2023visual} and CLIP \cite{radford2021learning}, while proficient with natural images, are not trained on scientific images. 
Developing a specialized LMM trained specifically on scientific images~\cite{li2024multimodal, wu2024scimmir} could potentially enhance its performance in interpreting scientific plots, figures, and tables, thereby improving its potential usage in scientific applications. 